\section{Stability under approximate recovery}
\label{sec:tpc40-approximate-recovery}

The exact Loewner argument has a direct biased version.  Let \(a\ne0\) be
arbitrary and define its gain vector \(g\in\C^R\) by
\begin{equation}
 g^T:=a^T\Phi.
 \label{eq:tpc40-approximate-gain}
\end{equation}
Thus exact unit-gain recovery is the special case \(g=\one_R\).  Retain
\(A=a^*W^{-1}a\), \(D=W^{1/2}\Phi\), \(b=W^{-1/2}a\), and
\(u=\overline b/\sqrt A\).

\begin{theorem}[Approximate weighted Loewner inequality]
\label{thm:tpc40-approximate-Loewner}
For every nonzero recovery coefficient \(a\),
\begin{equation}
 \boxed{
 G-\frac1A\overline g\,g^T
 =D^*(I_m-uu^*)D\succeq0.}
 \label{eq:tpc40-approximate-Loewner-decomposition}
\end{equation}
Consequently,
\begin{equation}
 A\lambda_{\max}(G)\ge\|g\|_2^2.
 \label{eq:tpc40-gain-norm-uncertainty}
\end{equation}
If, for some \(0\le\varepsilon\le1\),
\begin{equation}
 \|g-\one_R\|_2\le\varepsilon\sqrt R,
 \label{eq:tpc40-relative-L2-recovery-error}
\end{equation}
then
\begin{equation}
 \boxed{
 A\lambda_{\max}(G)
 \ge R(1-\varepsilon)^2.}
 \label{eq:tpc40-approximate-uncertainty-product}
\end{equation}
\end{theorem}

\begin{proof}
The definition of \(g\) gives
\[
 u^*D=\frac{b^TD}{\sqrt A}=\frac1{\sqrt A}g^T,
 \qquad
 D^*u=\frac1{\sqrt A}\overline g.
\]
This proves the factorization in
\eqref{eq:tpc40-approximate-Loewner-decomposition}.  The rank-one matrix
\(A^{-1}\overline g\,g^T\) has largest eigenvalue
\(A^{-1}\|g\|_2^2\), proving
\eqref{eq:tpc40-gain-norm-uncertainty}.  Finally, the reverse triangle
inequality and \eqref{eq:tpc40-relative-L2-recovery-error} give
\[
 \|g\|_2
 \ge\|\one_R\|_2-\|g-\one_R\|_2
 \ge(1-\varepsilon)\sqrt R.
\]
Squaring completes the proof.
\end{proof}

\begin{corollary}[Coordinatewise error]
\label{cor:tpc40-coordinatewise-recovery-error}
If
\begin{equation}
 \max_{1\le\alpha\le R}|g_\alpha-1|\le\varepsilon
 \qquad(0\le\varepsilon\le1),
 \label{eq:tpc40-coordinatewise-error-assumption}
\end{equation}
then \eqref{eq:tpc40-approximate-uncertainty-product} holds.
\end{corollary}

\begin{proof}
The coordinatewise assumption implies
\(\|g-\one_R\|_2\le\varepsilon\sqrt R\).
\end{proof}

The norm and normalization in
\eqref{eq:tpc40-relative-L2-recovery-error} are essential.  Without a
quantitative condition that forces \(\|g\|_2\) to remain close to
\(\sqrt R\), no lower bound with the factor
\((1-\varepsilon)^2\) follows.  Under the same error hypothesis, but for
unrestricted \(\varepsilon\ge0\), the proof yields the always-valid form
\begin{equation}
 A\lambda_{\max}(G)
 \ge R(1-\varepsilon)_+^2,
 \qquad (x)_+:=\max\{x,0\}.
 \label{eq:tpc40-approximate-positive-part}
\end{equation}
